\documentclass[11pt,letterpaper]{article}

\usepackage{preamble}

\title{Per-Person Human Action Recognition for Industrial VisionOps:\\
A Notebook-Driven Pipeline with V-JEPA~2 and InHARD}

\author{%
  Landy Haydee Schlebach Osorio \and
  Carlos Pano Hern{\'a}ndez \and
  Carlos Fernando Del Castillo Rey\\[0.6em]
  \small Tec de Monterrey, Campus Estado de M{\'e}xico --- Equipo MNA-V\\[0.3em]
  \small Academic advisor: Dr.~Gerardo Camacho\\
  \small Industrial sponsor: Alignity IQ Edge, LLC
}

\date{June 2026}

\begin{document}
\maketitle

\begin{abstract}
Manufacturing digital twins require timely interpretation of shop-floor video, yet scene-level human action recognition (\har{}) collapses distinct operator behaviors into a single label.
This manuscript documents a reproducible notebook pipeline for \emph{per-person} industrial \har{} built on the Industrial Human Action Recognition Dataset (\inhard{}; \num{5303} clips, fourteen meta-actions).
Frozen V-JEPA~2 embeddings (\texttt{facebook/vjepa2-vitl-fpc64-256}) feed a trainable multilayer perceptron head; YOLOv8 and ByteTrack maintain per-track crop buffers for live inference.
Exploratory analysis confirms severe class imbalance (e.g., \num{1378} \emph{Assemble system} clips versus \num{39} \emph{Take subsystem} clips).
A pilot configuration with five stratified clips per class ($n=\num{70}$) achieved \num{21.4}\% holdout accuracy and macro F1 \num{0.123}, with near-uniform confidence on mock-video sessions ($\approx\num{7.7}\%$ mean).
A scaled configuration targeting one hundred clips per class yields \num{1266} available samples under current disk inventory.
Versioned session logs, configuration fingerprinting, and automated analysis notebooks (steps~00--06) support iterative model comparison.
Results indicate that architectural integration is operational while discriminative performance remains sample-limited; completion of the scaled training run constitutes the immediate research priority.
\end{abstract}

\section{Introduction}
\label{sec:introduction}

Manufacturing supervisors often lack continuous visibility into shop-floor activity.
Cameras are widely deployed, yet raw video seldom translates into timely operational understanding.
Delays between an incorrect action and corrective feedback can extend to entire shifts, which affects productivity, quality, and training effectiveness.
Industry~4.0 initiatives therefore require systems that interpret visual streams while operations unfold rather than after reports are compiled.

\visionops{} denotes an industrial visual-intelligence prototype developed under academic--industry collaboration (Alignity IQ Edge, LLC; Tec de Monterrey, Campus Estado de M{\'e}xico).
The prototype follows three operational pillars---observe, guide, and improve---and targets a \emph{digital twin of the plant} derived from multiple synchronized video sources.
Human action recognition (\har{}) constitutes a central perception module: classifying operator meta-actions such as component handling, tool use, and sheet consultation.

A recurring failure mode in multi-operator scenes is \emph{scene-level} classification, in which a single label is assigned to an entire frame although several persons perform distinct actions simultaneously.
The present study addresses this limitation through a \emph{per-person} pipeline: person detection, multi-object tracking, per-track crop buffering, frozen video embeddings, and a lightweight classifier head.
Experiments use the Industrial Human Action Recognition Dataset (\inhard{})~\cite{gaspar2020inhard}, which provides segmented RGB clips for fourteen meta-actions recorded during collaborative assembly.

The contributions of the current iteration are threefold.
First, a reproducible notebook pipeline (steps~00--06) documents exploratory data analysis, embedding extraction, classifier training, live evaluation, session logging, and quantitative model analysis.
Second, stratified sampling and configuration fingerprinting enable controlled scaling experiments (e.g., five versus one hundred clips per class) without redundant re-computation.
Third, versioned session logs capture model identity, track identifiers, confidence scores, frame snapshots, and embedding artifacts to support iteration-over-iteration comparison.

The remainder of the manuscript is organized as follows.
\Cref{sec:literature} reviews related work on industrial \har{} and video foundation models.
\Cref{sec:methods} describes the dataset, pipeline architecture, and evaluation protocol.
\Cref{sec:results} reports exploratory statistics and classifier metrics from completed runs.
\Cref{sec:discussion} interprets findings and limitations.
\Cref{sec:conclusion} summarizes conclusions and planned extensions.

\section{Literature Review}
\label{sec:literature}

\subsection{Industrial human action recognition}

Industrial \har{} differs from general activity recognition because actions are short, repetitive, and constrained by workstation layout.
Gaspar et al.\ introduced \inhard{} as a segmented-video benchmark for collaborative assembly, annotating meta-actions at clip level rather than frame level~\cite{gaspar2020inhard}.
The dataset has become a reference for comparing embedding-based classifiers under class imbalance and blocked actions (e.g., full-system assembly segments that dominate volume).

Prior industrial prototypes often aggregate predictions at scene level, which simplifies deployment but collapses distinct operator behaviors into one label.
Multi-person tracking combined with crop-level inference offers an alternative aligned with supervisory use cases where accountability is assigned per operator.

\subsection{Video foundation models and frozen embeddings}

Self-supervised video encoders reduce the need for large labeled corpora when only a linear or multilayer perceptron (\textsc{mlp}) head must be trained.
Joint-Embedding Predictive Architectures (\textsc{jepa}) learn latent dynamics from unlabeled video~\cite{assran2025vjepa2}.
The publicly released V-JEPA~2 ViT-L/64 model provides fixed-length embeddings from short clip tensors and has been adopted in recent industrial \har{} baselines alongside image encoders such as DINOv2~\cite{ochs2023dinov2}.

Motion-contrastive extensions (MC-JEPA) further specialize temporal representations~\cite{wang2023mcjepa}.
The present iteration focuses on a frozen V-JEPA~2 backbone with a trainable \textsc{mlp} head to isolate the effect of sampling strategy and per-person inference before comparing alternative backbones.

\subsection{Detection, tracking, and operational logging}

Person detection commonly relies on single-stage detectors such as YOLOv8~\cite{jocher2023yolov8}.
Association-based trackers (e.g., ByteTrack) maintain identity-stable trajectories across frames~\cite{zhang2022bytetrack}.
For operational audit trails, structured logs that bind model version, subject track, predicted probabilities, and visual evidence support continuous improvement cycles analogous to experiment tracking in machine learning systems.

\subsection{Research gap}

Existing demonstrations of industrial \har{} rarely publish end-to-end notebook pipelines that connect exploratory analysis, scaled training schedules, per-person live inference, and session-level telemetry within one version-controlled repository.
The current work fills this gap for the \visionops{} research track by documenting procedures and preliminary metrics at multiple sampling scales.

\section{Methods and Procedures}
\label{sec:methods}

\subsection{Dataset and exploratory analysis}

The full \inhard{} corpus comprises \num{5303} segmented \texttt{.mp4} clips across fourteen meta-action folders under \texttt{Segmented/RGBSegmented/}.
Exploratory analysis (notebooks~01 and~01b) verified CSV--disk consistency, computed per-class counts, clip duration statistics (mean duration \SI{2.42}{\second}, standard deviation \SI{1.77}{\second}), and session breakdowns by subject and operation.
Twelve classes are designated \emph{trainable} industrial actions; two classes (\emph{Assemble system}, \emph{No action}) are often excluded in production configurations because of volume dominance or ambiguous semantics, although the current experiments include all fourteen classes unless noted.

Class imbalance is substantial: \emph{Assemble system} contains \num{1378} clips whereas \emph{Take subsystem} contains \num{39}.
Stratified sampling draws up to $k$ clips per class with fixed random seed (\num{42}) to construct controlled training subsets.

\subsection{Notebook pipeline architecture}

Table~\ref{tab:pipeline} summarizes the notebook sequence orchestrated from \texttt{00\_Pipeline\_Run\_All.ipynb}.
A companion notebook (\texttt{00\_b\_Pipeline\_Resume.ipynb}) skips embedding extraction and classifier training when configuration fingerprints match stored manifests (\texttt{embeddings.manifest.json}, \texttt{*.pipeline.json}).

\begin{table}[htbp]
\centering
\caption{Notebook pipeline stages and primary artifacts.}
\label{tab:pipeline}
\begin{tabularx}{\linewidth}{@{}l X l@{}}
\toprule
Step & Function & Output \\
\midrule
01 / 01b & Data strategy and \inhard{} EDA & \texttt{outputs/inhard\_eda/} \\
02 & V-JEPA~2 embedding extraction & \texttt{embeddings.npz} \\
03 & \textsc{mlp} classifier training & \texttt{checkpoints/*.pt} \\
04 & Per-person eval video & \texttt{perperson\_eval.mp4} \\
05 & Live OpenCV demo & \texttt{lib/live\_app.py} \\
06 & Metrics and session analysis & \texttt{outputs/har\_analysis/} \\
\bottomrule
\end{tabularx}
\end{table}

\subsection{Embedding extraction and classifier training}

Each training clip is read with OpenCV, temporally subsampled to sixteen RGB frames at \SI{224}{px} resolution, normalized with ImageNet statistics, and encoded with the Hugging Face model \texttt{facebook/vjepa2-vitl-fpc64-256}.
Embeddings are mean-pooled over tokens and $\ell_2$-normalized.
A three-layer \textsc{mlp} (hidden width \num{512}, dropout \num{0.3}) is trained with AdamW ($\text{lr}=\num{e-3}$, weight decay \num{e-4}) for twenty-five epochs using an \num{80}/\num{20} stratified holdout split.

Pilot configuration \texttt{all14\_5each} uses five clips per class ($n=\num{70}$).
Scaled configuration \texttt{all14\_100each} targets one hundred clips per class (up to \num{1400} embeddings); available disk counts cap four rare classes below one hundred, yielding $n=\num{1266}$ in the verified sample manifest.

\subsection{Per-person inference stack}

Live and eval modes process mock industrial video (\texttt{Industrial-One.mp4}) or webcam input.
YOLOv8n detects persons; ByteTrack assigns \texttt{track\_id} values.
Each track maintains a deque of cropped person regions (default buffer: thirty-two frames).
Every sixteen frames, when the buffer exceeds half capacity, crops are embedded and classified.
A dwell smoother requires two consecutive agreeing predictions before the displayed label changes, reducing flicker.

The \har{} session logger writes JSONL events under \texttt{outputs/har\_sessions/YYYY-MM-DD/}, storing model tag, classifier version, all class probabilities, confidence, inference latency, and---on confirmed detections---annotated frame JPEGs and \texttt{.npy} embedding files.
A global \texttt{index.csv} supports cross-run comparison.

\subsection{Evaluation metrics}

Holdout evaluation reports accuracy, macro and weighted F1, and per-class precision/recall from \texttt{06\_Model\_and\_Session\_Analysis.ipynb}.
Live-session analysis aggregates confidence histograms, label frequencies, and latency distributions.
Ground-truth labels are unavailable for mock-video live inference; therefore live metrics are descriptive rather than supervised.

\section{Results}
\label{sec:results}

\subsection{Exploratory data analysis}

Figure~\ref{fig:eda-balance} shows the long-tailed class distribution across all \num{5303} clips.
Trainable actions (twelve classes, \num{3425} clips) exhibit lower mean duration (\SI{1.85}{\second}) than the full corpus because blocked classes include longer assembly segments.

\begin{figure}[htbp]
\centering
\includegraphics[width=0.92\linewidth]{01_class_balance_all.png}
\caption{Per-class clip counts in the full \inhard{} corpus ($n=\num{5303}$).}
\label{fig:eda-balance}
\end{figure}

\subsection{Pilot classifier (\texttt{all14\_5each})}

Table~\ref{tab:pilot-metrics} summarizes holdout metrics for the pilot run trained on seventy stratified clips (five per class, seed~\num{42}).

\begin{table}[htbp]
\centering
\caption{Holdout metrics for checkpoint \texttt{har\_vjepa\_all14\_5each.pt} ($n_{\text{train}}=\num{56}$, $n_{\text{holdout}}=\num{14}$).}
\label{tab:pilot-metrics}
\begin{tabular}{lr}
\toprule
Metric & Value \\
\midrule
Accuracy & \num{21.4}\% \\
Macro F1 & \num{0.123} \\
Weighted F1 & \num{0.123} \\
Macro precision & \num{0.104} \\
Macro recall & \num{0.214} \\
\bottomrule
\end{tabular}
\end{table}

Figure~\ref{fig:confusion} presents the row-normalized confusion matrix on the holdout split.
Figure~\ref{fig:prf1} reports per-class precision, recall, and F1.
With only one holdout sample per class on average, several classes register zero precision or recall; the loss curve (Figure~\ref{fig:loss}) shows minimal separation between training and validation cross-entropy ($\approx\num{2.63}$), consistent with near-chance performance at this sample size.

\begin{figure}[htbp]
\centering
\includegraphics[width=0.85\linewidth]{02_confusion_matrix.png}
\caption{Row-normalized confusion matrix (pilot holdout).}
\label{fig:confusion}
\end{figure}

\begin{figure}[htbp]
\centering
\includegraphics[width=\linewidth]{03_per_class_prf1.png}
\caption{Per-class precision, recall, and F1 on the pilot holdout split.}
\label{fig:prf1}
\end{figure}

\begin{figure}[htbp]
\centering
\includegraphics[width=0.72\linewidth]{04_training_loss.png}
\caption{Training and validation loss across twenty-five epochs (pilot run).}
\label{fig:loss}
\end{figure}

\subsection{Per-person live and eval sessions}

Five logged sessions (eval and live) produced \num{93} inference events and \num{38} confirmed detections for model tag \texttt{all14\_5each}.
Mean predicted confidence was \num{7.7}\%; mean end-to-end latency was approximately \SI{3671}{\milli\second} on CPU (V-JEPA~2 embedding plus \textsc{mlp} forward pass per track).

Figure~\ref{fig:confidence} illustrates the confidence distribution.
Figure~\ref{fig:labels} shows that \emph{Put down component} dominated predicted labels on mock video, suggesting either scene bias or insufficient discriminative training signal at the pilot scale.

\begin{figure}[htbp]
\centering
\includegraphics[width=\linewidth]{05_session_confidence.png}
\caption{Confidence distributions for all inferences and confirmed detections.}
\label{fig:confidence}
\end{figure}

\begin{figure}[htbp]
\centering
\includegraphics[width=0.88\linewidth]{06_session_label_distribution.png}
\caption{Predicted action frequencies across logged live/eval sessions (pilot model).}
\label{fig:labels}
\end{figure}

The per-person eval video (\num{102} frames processed) demonstrated simultaneous tracks with distinct overlay labels, validating the architectural separation of inference streams by \texttt{track\_id}.

\subsection{Scaled training iteration (\texttt{all14\_100each})}

Configuration scaling to one hundred clips per class was implemented with fingerprint-based resume (\texttt{00\_b}).
Verified stratified sampling yields \num{1266} clips because four meta-actions contain fewer than one hundred segments on disk (\emph{Take subsystem}: \num{39}; \emph{Put down measuring rod}: \num{74}; \emph{Take measuring rod}: \num{76}; \emph{Put down subsystem}: \num{77}).
Full embedding extraction for the scaled run was in progress at the time of writing; holdout metrics for \texttt{all14\_100each} will be reported upon completion in \texttt{outputs/har\_analysis/}.

\section{Discussion}
\label{sec:discussion}

\subsection{Interpretation of pilot metrics}

Holdout accuracy of \num{21.4}\% on fourteen classes exceeds uniform guessing (\num{7.1}\%) yet remains far from operational deployment thresholds.
The macro F1 score of \num{0.123} indicates that precision and recall are uneven across meta-actions, which is expected when only five labeled clips per class constrain the \textsc{mlp} head.
Softmax outputs clustered near uniform values (typical top confidence $\approx\num{7}\%$--\num{8}\%) further suggest that the classifier has not learned sharp decision boundaries at the pilot scale.

These outcomes do not invalidate the pipeline; rather, they quantify the sample complexity required before per-person live guidance becomes reliable.
The exploratory analysis confirms ample unused data (\num{5303} clips), so performance gains likely depend on labeled exposure rather than architectural changes alone.

\subsection{Per-person inference versus scene-level baselines}

Assigning labels per \texttt{track\_id} avoids propagating one scene-level prediction to every visible operator.
Session logs demonstrate multiple concurrent tracks with independent label histories.
Nevertheless, crop quality, occlusion, and identity switches remain error sources that were not fully quantified in the absence of frame-level ground truth on mock video.

\subsection{Operational logging and reproducibility}

Versioned session directories bind each detection to model tag (\texttt{all14\_5each}), embedding model identifier, checkpoint name, and timestamp.
This structure supports comparative studies as sampling scales increase.
Configuration fingerprints prevent accidental reuse of mismatched embeddings and checkpoints---a practical requirement when iteration cycles span hours of CPU extraction.

\subsection{Limitations}

Several limitations should be acknowledged openly.
First, experiments used CPU inference; latency figures may decrease on GPU or edge accelerators but were not measured here.
Second, live evaluation lacked per-frame ground truth, precluding supervised session-level F1.
Third, rare \inhard{} classes cap stratified sampling below the nominal one-hundred-clips target.
Fourth, blocked production classes (\emph{Assemble system}, \emph{No action}) may require exclusion or hierarchical modeling to reflect deployment policy.
Fifth, the manuscript reports completed pilot metrics; scaled \texttt{all14\_100each} results remain pending.

\subsection{Future work}

Subsequent iterations should complete scaled training runs, report holdout and live-session metrics in `har-research/04_Analysis_and_Visualization.ipynb`, compare frozen V-JEPA~2 against DINOv2 heads under identical sampling, and evaluate on held-out \inhard{} subjects not seen during stratified sampling.
The production demo integrates per-person inference through the unified VisionOps FastAPI backend (\texttt{vision-ops-backend}, port~8000) with live UI at \texttt{/live} and human review at \texttt{/har-hitl}.

\section{Conclusion}
\label{sec:conclusion}

The present study documented a notebook-driven, per-person human action recognition pipeline for industrial \visionops{} prototypes.
Exploratory analysis characterized \num{5303} \inhard{} clips with pronounced class imbalance.
A frozen V-JEPA~2 encoder with a trainable \textsc{mlp} head was integrated with YOLOv8 detection, ByteTrack association, dwell-smoothed labeling, and versioned session telemetry.

Pilot experiments at five clips per class ($n=\num{70}$) established end-to-end reproducibility but yielded modest holdout accuracy (\num{21.4}\%, macro F1 \num{0.123}), indicating that additional labeled exposure is necessary before operator-facing guidance can be trusted.
Scaling to approximately one hundred clips per class ($n=\num{1266} available) with resume-safe caching represents the active iteration path.

The pipeline, analysis notebooks, and session logs collectively provide an auditable foundation for iterative improvement aligned with the observe--guide--improve objectives of digital-twin-enabled manufacturing supervision.


\bibliographystyle{plainnat}
\bibliography{references}

\end{document}
